\documentclass[11pt,a4paper]{article}
\usepackage{gtu-el-yazisi}
\begin{document}
\notbilgi{Makine Öğrenmesi}{NLP: Tokenization, TF-IDF, Word2Vec}{Ders-10.pdf}

\sayfa{1}
\section{NLP (Natural language processing)}
\subsection{Tokenization}

\underline{My}\textsubscript{\scriptsize token} \underline{name} \underline{is} \underline{Mehmed}\textcircled{.}\textcolor{kirmizikalem}{\,?}

Morfolojik gitmek istersek (şekle göre)

Benim adım Mehmet.

\textcolor{turuncukalem}{Ben-im ad-ım}

$\uparrow$ kelimelere bölerek word piece

$\rightarrow$ Morfolojik gramere göre

Türkçede 10000 kelime

\quad 10 000

\quad ahmed $\to$ 2000
\[ \text{\small one hot encoding}\qquad \begin{bmatrix}0\\0\\0\\0\\1\\0\\0\end{bmatrix} \]

Elimizde dokümanlar var

\elcizimi[0.45]{sekiller/Ders-10-s1-1.png}[1]

TF\\
{\small Term frequency}

Tüm dokümanlarda iki kelime aynı frekansla geçmemeli

\underline{TF-IDF} $\to$ baskınlığı azaltmak için

\section{Word2Vec}
\subsection{CBOW}
Continuous Bag of Words

% (şekil içindeki not: Embedding vektörü, kaçlık vektör istiyorsan yapabilirim)
\elcizimi[0.9]{sekiller/Ders-10-s1-2.png}[1]

\sayfa{2}
\section{Skip gram / n-gram}
\textcolor{kirmizikalem}{2 gram iki komşuluk}

tersten de yapabilirim

\elcizimi[0.7]{sekiller/Ders-10-s2-1.png}[2]

\section{Global vectors for word Representation (Glove)}
\[
J=\sum_{i,j} f(x_{ij})\Bigl(\underbrace{w_i^T}_{\substack{\text{\scriptsize bunu bulmaya}\\ \text{\scriptsize çalıştığımız}\\ \text{\scriptsize kelimenin}\\ \text{\scriptsize ağırlığı}}}\,\widetilde{w}_j+b_i+\widetilde{b}_j-\underbrace{\log X_{ij}}_{\substack{\text{\scriptsize i'ninci doküman}\\ \text{\scriptsize j'inci kelimenin}\\ \text{\scriptsize geçme frekansı}}}\Bigr)^2
\]

\medskip
{\color{kirmizikalem}\noindent\rule{0.55\linewidth}{0.6pt}}
\medskip

% (şekil içindeki not: apple = t1+t2; app/le, t1 t2, subword; word embedding; Fast Text)
\elcizimi[0.45]{sekiller/Ders-10-s2-2.png}[2]

\section{Static embedding}
\textcolor{kirmizikalem}{Covid geldi yeni kelime}

\section{Benim adım Mehmet / Seq 2Seq}
\textcolor{kirmizikalem}{$\swarrow$ Supervised}

\elcizimi[0.5]{sekiller/Ders-10-s2-3.png}[2]

encoder

\textcolor{kirmizikalem}{\underline{RNN}} $\to$ exploding vanishing \qquad unutma fazla \okunamadi{hatırlama}

\begin{tabular}{c@{\quad}c@{\quad}c@{\quad}c}
\textcolor{kirmizikalem}{My} & \textcolor{kirmizikalem}{name} & \textcolor{kirmizikalem}{is} & \textcolor{kirmizikalem}{Mehmet}\\
$\widetilde{w}_1$ & $\widetilde{w}_2$ & $\widetilde{w}_3$ & $\widetilde{w}_4$
\end{tabular}
\qquad sonraki güncellemede LSTM

Bu metin kısa olmadıktan sonra context'i yakalayamıyor.

\textcolor{kirmizikalem}{\underline{decoder}}

\elcizimi[0.4]{sekiller/Ders-10-s2-4.png}[2]

\sayfa{3}
\section{Attention Mechanism}
% (şekil içindeki not: Context, metni temsil edecek gibi, s_1 h)
\elcizimi[0.7]{sekiller/Ders-10-s3-1.png}[3]

kelimeleri hidden layerle çarpıyoruz

softmax $\downarrow$ \quad $w\cdot h$
\[ [\alpha_1,\alpha_2,\dots,\alpha_{100}] \]

hangi kelime ne kadar önemli

sonra dönmek yerine MLP koyuyoruz

\elcizimi[0.5]{sekiller/Ders-10-s3-2.png}[3]

paralelde çalışabiliyor

\textcolor{kirmizikalem}{\underline{Encoding:}} \hspace{5em} \textcolor{kirmizikalem}{\underline{Decoder:}} \hspace{2em} \textcolor{kirmizikalem}{GPT sadece decoder}

% (şekil içindeki notlar: input encoding; output input embedding; Dynamic Embedding)
\elcizimi[0.85]{sekiller/Ders-10-s3-3.png}[3]

\fbox{\begin{minipage}{0.42\linewidth}
ali kopu at

a/li/ /ko/pu/ /at/

arada hataları düzeltmek için token \okunamadi{sığıyor}
\end{minipage}}

\medskip
\textsuperscript{2019}\underline{BERT}: Sadece Encoder , kelimeleri ayırabilirsin

\hspace*{1.5em}BERT NLP'nin YOLO'su

BERT : Bidirectional Encoder Representations from Transformers

\sayfa{4}
\underline{ELMO}: Embeddings from language model

context mekanizması yok

\elcizimi[0.5]{sekiller/Ders-10-s4-1.png}[4]

\end{document}
